\documentclass[11pt,a4paper]{article}

\usepackage{iftex}
\ifPDFTeX
  \usepackage[utf8]{inputenc}
  \usepackage[T1]{fontenc}
  \usepackage{lmodern}
\else
  \usepackage{fontspec}
  \setmainfont{Times New Roman}
\fi
\usepackage[magyar]{babel}
\usepackage{microtype}
\usepackage[a4paper,margin=25mm,headheight=14pt]{geometry}
\usepackage{amsmath,amssymb,mathtools}
\usepackage{booktabs,longtable,array,tabularx}
\usepackage{enumitem}
\usepackage{xcolor}
\usepackage{fancyhdr}
\usepackage{hyperref}
\usepackage{url}
\providecommand*{\theHpage}{\thepage}

\definecolor{navy}{HTML}{17324D}
\definecolor{teal}{HTML}{087E8B}
\definecolor{lightblue}{HTML}{EDF4F8}
\hypersetup{
  colorlinks=true,
  linkcolor=navy,
  urlcolor=teal,
  citecolor=navy,
  pdftitle={DC motor elektromechanikai modelljének PINN-alapú vizsgálata},
  pdfauthor={Projektterv}
}
\setlength{\parindent}{0pt}
\setlength{\parskip}{5pt}
\setlist{itemsep=2pt,topsep=3pt}
\setlength{\emergencystretch}{2em}
\renewcommand{\arraystretch}{1.15}
\pagestyle{fancy}
\fancyhf{}
\fancyhead[L]{DC motor PINN projekt}
\fancyhead[R]{Műszaki dokumentáció}
\fancyfoot[C]{\thepage}
\renewcommand{\headrulewidth}{0.4pt}

\newcommand{\dd}{\mathrm{d}}
\newcommand{\vect}[1]{\boldsymbol{#1}}
\newcommand{\unit}[1]{\,\mathrm{#1}}
\newcommand{\RMSE}{\mathrm{RMSE}}
\newcommand{\MAE}{\mathrm{MAE}}

\begin{document}

\hypersetup{pageanchor=false}
\begin{titlepage}
\centering
\vspace*{18mm}
{\Large\color{teal} PROJEKTDOKUMENTÁCIÓ\par}
\vspace{12mm}
{\huge\bfseries\color{navy} Egyenáramú motor\\[2mm] elektromechanikai modelljének\\[2mm] PINN-alapú vizsgálata\par}
\vspace{9mm}
{\Large Fizikai modell, állapotbecslés és rendszerazonosítás\par}
\vspace{20mm}
\begin{minipage}{0.83\textwidth}
\large
A dokumentum egy állandómágneses, kefés DC motor fokozatos PINN-vizsgálatának mérnöki tervét rögzíti. A központi kérdés az, hogy a mérési adatok és a fizikai egyenletek együttes használata miként segítheti az állapotok, a motorparaméterek és az ismeretlen terhelőnyomaték becslését.
\end{minipage}
\vfill
{\large 2026. november 3.\par}
\vspace{5mm}
\end{titlepage}

\hypersetup{pageanchor=true}
\pagenumbering{roman}
\renewcommand*{\theHpage}{roman.\roman{page}}
\tableofcontents
\newpage
\pagenumbering{arabic}
\renewcommand*{\theHpage}{arabic.\arabic{page}}

\section{Összefoglaló és célkitűzés}

A projekt célja egy időfüggő, elektromechanikailag csatolt rendszer vizsgálata fizikai információt felhasználó neurális hálózattal (physics-informed neural network, PINN). A kiinduló rendszer egy állandómágneses, kefés DC motor. A modell az armatúraáram, a szögsebesség és a tengelyhelyzet dinamikáját kapcsolja össze a kapocsfeszültséggel és a külső terhelőnyomatékkal.

A fejlesztés fokozatos: először ismert paraméterekkel készül előrejelző modell, majd következik a paraméterazonosítás, a ritka mérésekből történő állapotbecslés, az ismeretlen terhelés rekonstruálása, végül a zajjal, modellhibával és korábban nem látott gerjesztésekkel szembeni robusztusság vizsgálata. A referenciaadatok elsődleges forrása MATLAB/Simulink szimuláció, amely kontrollált igazságadatot ad. Ezt egészítheti ki egy nyilvános, kísérleti DC motor adatbázis; az adatbázisban ténylegesen elérhető jeleket és mértékegységeket a letöltött fájlok alapján kell véglegesíteni.

A projekt eredménye nem pusztán egy előrejelző hálózat, hanem összehasonlítható kísérletsorozat: megmutatja, hogy a fizikai veszteségtag mikor javítja az állapotbecslést és az általánosítást, mely paraméterek azonosíthatók a választott gerjesztésekkel, valamint hol jelentkeznek a modell és a mérési adatok korlátai.

\section{Hatókör, rendszerhatár és jelölések}

\subsection{Vizsgált rendszer}

A rendszerhatár a motor armatúrakörét, a forgórészt, a tengelyre redukált tehetetlenséget és az ezekre ható külső terhelést tartalmazza. A bemenet az armatúra kapocsfeszültsége. Az alapállapot-vektor:
\begin{equation}
  \vect{x}(t)=
  \begin{bmatrix}
    i(t) & \omega(t) & \theta(t)
  \end{bmatrix}^{\mathsf T},
  \label{eq:state}
\end{equation}
ahol az állapotok rendre az armatúraáram, a szögsebesség és a tengely szöghelyzete. Az alapbemenet $u_a(t)$, az ismeretlen vagy ismert zavaró bemenet pedig a terhelőnyomaték $M_L(t)$.

\subsection{Jelölésrendszer}

\small
\renewcommand{\arraystretch}{1.0}
\setlength{\parskip}{0pt}
\begin{longtable}{>{\raggedright\arraybackslash}p{0.23\textwidth} >{\raggedright\arraybackslash}p{0.70\textwidth}}
\toprule
Jelölés & Jelentés és SI-mértékegység \\
\midrule
\endhead
$t$ & Idő, $\unit{s}$.\\
$u_a$ & Armatúra kapocsfeszültsége, $\unit{V}$.\\
$i$ & Armatúraáram, $\unit{A}$.\\
$R$ & Armatúraellenállás, $\unit{\Omega}$.\\
$L$ & Armatúrainduktivitás, $\unit{H}$.\\
$e_b$ & Visszaindukált feszültség, $\unit{V}$.\\
$k_e$ & Visszaindukált feszültség állandója, $\unit{V\,s\,rad^{-1}}$.\\
$k_t$ & Nyomatékállandó, $\unit{N\,m\,A^{-1}}$.\\
$M_m$ & Elektromágneses motor-nyomaték, $\unit{N\,m}$.\\
$M_L$ & Külső terhelőnyomaték, $\unit{N\,m}$.\\
$J$ & Egyenértékű tengely-tehetetlenségi nyomaték, $\unit{kg\,m^2}$.\\
$b$ & Viszkózus súrlódási tényező, $\unit{N\,m\,s\,rad^{-1}}$.\\
$\omega$ & Szögsebesség, $\unit{rad\,s^{-1}}$.\\
$\theta$ & Szöghelyzet, $\unit{rad}$.\\
$\alpha=\dot\omega$ & Szöggyorsulás, $\unit{rad\,s^{-2}}$.\\
$\hat{z}$, $z$ & Becsült, illetve referenciaérték.\\
\bottomrule
\end{longtable}
\normalsize
\renewcommand{\arraystretch}{1.15}
\setlength{\parskip}{5pt}

A pont a teljes idő szerinti deriváltat jelöli, például $\dot{i}=\dd i/\dd t$. A pozitív forgásirányban a motor nyomatéka gyorsítja a rendszert, a pozitív $M_L$ pedig ezzel ellentétes, fékező nyomatékként szerepel. A szögsebességet radián per másodpercben kell használni; ha a forrás fordulatszámot ad meg, akkor
\begin{equation}
  \omega = \frac{2\pi}{60}n,
  \label{eq:rpm}
\end{equation}
ahol $n$ a percenkénti fordulatszám.

\section{A fizikai DC motor modellje}

\subsection{Armatúra feszültségegyenlete}

Az armatúrakört soros $R$--$L$ tag és a forgás miatt keletkező ellenfeszültség írja le. A Kirchhoff-féle huroktörvény szerint
\begin{equation}
  u_a(t)=R\,i(t)+L\frac{\dd i(t)}{\dd t}+e_b(t),
  \label{eq:armature}
\end{equation}
ahol az állandómágneses motor lineáris tartományában
\begin{equation}
  e_b(t)=k_e\,\omega(t).
  \label{eq:backemf}
\end{equation}
Ebből az áram állapotegyenlete:
\begin{equation}
  \dot{i}(t)=\frac{u_a(t)-R\,i(t)-k_e\omega(t)}{L}.
  \label{eq:idot}
\end{equation}
Az $Ri$ tag az ellenálláson eső feszültség, az $L\dot{i}$ az áramváltozáshoz szükséges induktív feszültség, a $k_e\omega$ pedig a forgás által indukált, a táplálással szemben ható feszültség.

\subsection{Elektromágneses nyomaték}

A motor nyomatéka a mágneses tér és az armatúraáram kölcsönhatásából származik. A lineáris modellben
\begin{equation}
  M_m(t)=k_t\,i(t).
  \label{eq:torque}
\end{equation}
Állandómágneses DC motor esetében, következetes SI-egységekkel, a teljesítményátalakítási kapcsolat $e_b i=M_m\omega$ alapján $k_e$ és $k_t$ számértéke azonos. Az összefüggés csak azonos definíció és SI-egységek mellett alkalmazható; adatlapokban a fordulatszám más egységben szerepelhet, ezért a konverziót ellenőrizni kell. A kezdeti vizsgálatban a két állandó külön becsülhető, később pedig az egyenlőség fizikailag indokolt kötésként is bevezethető.

\subsection{Mechanikai dinamika és pozíció}

A forgómozgás Newton-törvénye szerint az eredő nyomaték egyenlő a tehetetlenség és a szöggyorsulás szorzatával. Viszkózus súrlódást és külső terhelőnyomatékot figyelembe véve:
\begin{equation}
  J\dot{\omega}(t)=M_m(t)-b\omega(t)-M_L(t).
  \label{eq:mechanics}
\end{equation}
A sebesség és a pozíció kinematikai kapcsolata:
\begin{equation}
  \dot{\theta}(t)=\omega(t).
  \label{eq:position}
\end{equation}
Ezekből a mechanikai állapotegyenlet:
\begin{equation}
  \dot{\omega}(t)=\frac{k_t i(t)-b\omega(t)-M_L(t)}{J}.
  \label{eq:omegadot}
\end{equation}
A teljes, háromállapotú modell tehát
\begin{equation}
\boxed{
\begin{aligned}
  \dot{i} &= \frac{u_a-Ri-k_e\omega}{L},\\
  \dot{\omega} &= \frac{k_t i-b\omega-M_L}{J},\\
  \dot{\theta} &= \omega.
\end{aligned}}
\label{eq:fullmodel}
\end{equation}
A terhelés lehet ismert bemeneti jel, becsülendő ismeretlen vagy a kezdeti kísérletekben nulla. A pozíció egy integráló állapot; ha az adatforrás csak fordulatszámot közöl, a pozíció közvetlen mérési veszteségtagja elhagyható, de a kinematikai egyenlet továbbra is használható.

\subsection{Alapfeltevések és érvényességi tartomány}

Az alapmodell a következő egyszerűsítéseket teszi:

\begin{itemize}
  \item állandómágneses, kefés DC motor, közel állandó mágneses fluxussal;
  \item állandó $R$, $L$, $k_e$, $k_t$, $J$ és $b$ egy-egy vizsgált kísérleten belül;
  \item merev tengely, a motor és a terhelés tehetetlensége egyetlen $J$ paraméterbe redukálva;
  \item lineáris, sebességgel arányos viszkózus súrlódás;
  \item az armatúrafeszültség a tényleges motor kapcsain értendő, nem automatikusan a PWM parancs;
  \item kefefeszültség, telítődés, fogazási nyomatéklüktetés, holtjáték és hőmérsékletfüggés az első modellben elhanyagolva.
\end{itemize}

Ezek nem univerzális igazságok: a mérési adatokkal ellenőrizendő modellezési feltevések. Ha a kis sebességű eltérések jelentősek, a Coulomb-súrlódás például $M_c\operatorname{sgn}(\omega)$ taggal bővítheti az egyenletet; ezt csak az alapmodell validálása után célszerű bevezetni.

\section{A PINN feladatai}

A PINN egy folytonos időfüggvényt közelít a hálózattal, majd automatikus deriválással kiértékeli a differenciálegyenletek maradékát. A mért adatok a tényleges kimenetekhez kötik a becslést, a fizikai egyenletek pedig a mintavételi pontok közötti viselkedést is korlátozzák.

\subsection{Előreirányú feladat (forward problem)}

Ismert $R,L,k_e,k_t,J,b$, ismert $u_a(t)$, ismert vagy rögzített $M_L(t)$ és kezdeti állapot mellett a hálózat az
\[
  t \longmapsto
  \begin{bmatrix}\hat{i}(t)&\hat{\omega}(t)&\hat{\theta}(t)\end{bmatrix}^{\mathsf T}
\]
leképezést tanulja. Ez az első lépcső: a cél az ismert paraméterű referencia-ODE reprodukálása. Először sűrű szimulációs megfigyelésekkel érdemes összevetni egy hagyományos numerikus ODE-megoldót; csak ezután csökkentsük a mért pontok számát.

\subsection{Inverz feladat és paraméterazonosítás}

Az inverz feladatban a fizikai paraméterek egy részét az adatokból kell becsülni. Például:
\begin{equation}
  \vect{p}=
  \begin{bmatrix}R&L&k_e&k_t&J&b\end{bmatrix}^{\mathsf T}.
\end{equation}
A hálózati súlyok mellett a kiválasztott paraméterek is tanítható változók. Pozitivitásuk biztosítható például $p=\exp(q)$ parametrizálással. Nem cél minden paramétert minden kísérletből egyszerre becsülni: a paraméterek összefüggése és a gerjesztés információtartalma korlátozhatja az egyértelmű azonosíthatóságot. Kezdetben célszerű egyes paramétereket adatlapból vagy külön mérésből rögzíteni, és csak néhány, jól gerjesztett paramétert becsülni.

\subsection{Állapotbecslés ritka vagy zajos mérésekből}

A hálózat azokat az állapotokat is folytonosan becsülheti, amelyeket csak ritkán vagy zajosan mérünk. Például mérhető lehet a fordulatszám és az áram, miközben a tanítás során csak ezek egy részét használjuk fel; a fizikai modell segíthet az időközi állapot rekonstruálásában. A pozíció becsülhető a $\dot{\theta}=\omega$ kapcsolatból, de az integrálási drift miatt a kezdeti pozíció rögzítése vagy időnkénti pozíciómérés fontos lehet.

Ez a feladat nem azonos az extrapolációval: az időtartományon belüli ritka minták közötti interpoláció általában könnyebb, mint korábban nem látott bemeneti profil vagy üzemi tartomány előrejelzése.

\subsection{Ismeretlen terhelőnyomaték becslése}

A mechanikai maradékból az ismeretlen terhelés formálisan kifejezhető:
\begin{equation}
  M_L(t)=k_t i(t)-b\omega(t)-J\dot{\omega}(t).
  \label{eq:loadformal}
\end{equation}
A közvetlen differenciálás zajérzékeny, ezért a PINN inkább együtt tanulja $\omega(t)$ sima reprezentációját és $M_L(t)$-t. A terhelés közelíthető külön neurális kimenettel, simító szplájnnal vagy intervallumonként állandó értékkel. A választás a terhelés fizikai jellegétől függjön: lépcsős terhelésre nem célszerű túl erős simasági büntetést alkalmazni.

A terhelés, a tehetetlenség és a súrlódás egyidejű becslése könnyen több, hasonlóan jó magyarázatot eredményezhet. Ezért a szimulációban ismert terhelési profilokkal kell demonstrálni a módszert, a valós adatokon pedig csak azonosítható paraméterkészletet szabad becsültnek tekinteni.

\section{Adatforrás-stratégia}

\subsection{Nyilvános kísérleti adatbázis}

A beszélgetésben kiválasztott kiinduló forrás a Mendeley Data OpenMCT DC motor adatkészlete. A jelenlegi, 2. verziójú rekord leírása 14 nyers adatfelvételt, több mint 60\,000 mintasort, motorsebességet, referenciajelet, szabályozási periódust, áramérzékelő nyers és szűrt jeleit, PWM működtetőjelet, valamint egyes kísérletekben szinkronizált digitális multiméteres árammérést sorol fel. A leírás APRBS rendszerazonosító gerjesztéseket, több mintavételi időt, PI-szabályozási méréseket, chirp gerjesztést, MATLAB szkripteket és feldolgozott MAT fájlokat is említ. Az adathalmaz CC BY 4.0 licencet jelez \cite{mendeley}.

A leírás alapján az adatkészlet hasznos lehet fordulatszám- és áramállapot-becsléshez, rendszerazonosítási gerjesztések összehasonlításához és valós méréseken végzett robusztussági próbához. A nyilvános leírás azonban nem igazolja, hogy minden futamban közvetlenül elérhető a motor kapocsfeszültsége, a tengely szöghelyzete vagy a külső terhelőnyomaték. Ezért a letöltés után kötelező az oszlopok, egységek, mintavételi idők, kalibráció és fájlverzió ellenőrzése.

Különösen fontos különbséget tenni a PWM parancs és a tényleges $u_a(t)$ kapocsfeszültség között. Ha csak PWM kitöltési tényező áll rendelkezésre, az armatúrafeszültséghez ismerni kell többek között a DC sín feszültségét és a meghajtó veszteségeit. Ezek hiányában a feszültség csak becsült bemenetként kezelhető, és az ebből számított paraméterek bizonytalanságát jelezni kell.

\subsection{MATLAB/Simulink szimuláció}

A Simulink-modell elsődleges célja ellenőrizhető referenciaadat és igazságértékek előállítása. A szimulációban minden állapot, paraméter és terhelés ismert, ezért külön mérhető a hálózat előrejelzési és azonosítási hibája. Javasolt gerjesztések:

\begin{itemize}
  \item lépcsős és többlépcsős feszültségprofilok az indítási és állandósult válasz vizsgálatára;
  \item APRBS vagy PRBS gerjesztés a szélesebb frekvenciatartományú azonosításhoz;
  \item szinuszos vagy chirp feszültségprofil;
  \item ismert lépcsős vagy szakaszonként változó terhelőnyomaték;
  \item zajjal, ritkább mintavételezéssel és szenzoreltolással módosított mérési változatok.
\end{itemize}

A gerjesztések amplitúdóját a kiválasztott motor megengedett üzemi tartományához kell igazítani. Minden adatsorhoz együtt kell tárolni a paramétereket, kezdeti állapotokat, bemeneti profilokat, solver-beállításokat, mintavételi időt és az esetleges zajmodell magját.

\subsection{Hibrid adatfolyam}

A javasolt sorrend: (1) PINN és numerikus megoldó ellenőrzése szimulációval; (2) az adatbázis jeleinek auditja és adatelőkészítés; (3) külön mérési futamokon validált tanítás; (4) szimuláció és mérés eredményeinek összevetése; (5) a szimulált modell eltérésének tudatos vizsgálata. A szimulált és valós adatokat nem szabad automatikusan egyetlen tanítóhalmazba összekeverni: a modell- és érzékelőkülönbségeket előbb külön kell mérni.

Az adatok felosztása lehetőleg teljes kísérleti futamok szerint történjen. Véletlen mintaszintű felosztás esetén az egymáshoz közeli időpontok könnyen egyszerre kerülnek tanító és teszt adathalmazba, ami túl optimista eredményt adhat.

\section{PINN architektúra és veszteségfüggvény}

\subsection{Hálózati felépítés}

Egy kezdeti, MATLAB-ban automatikus deriválással tanítható architektúra teljesen összekapcsolt, sima aktivációjú (például $\tanh$) hálózat. Az egyetlen kísérlet időfüggvényeit a hálózat:
\begin{equation}
  f_{\vect{\phi}}(t)\longrightarrow
  \begin{bmatrix}
    \hat{i}(t)&\hat{\omega}(t)&\hat{\theta}(t)
  \end{bmatrix}^{\mathsf T}
  \label{eq:network}
\end{equation}
alakban közelíti, ahol $\vect{\phi}$ a tanítható hálózati paraméterek halmaza. A tanulás során automatikus deriválással előáll:
$\dot{\hat{i}}$, $\dot{\hat{\omega}}$ és $\dot{\hat{\theta}}$.

Egyetlen ismert gerjesztésű futamnál a bemenet lehet a normalizált idő; a mért vagy előírt $u_a(t)$ és $M_L(t)$ a fizikai maradék számításakor, a kollokációs pontokban kerül kiértékelésre. Több eltérő gerjesztés közös tanításánál a hálózatnak a futamot vagy gerjesztést is meg kell különböztetnie: ehhez használható kísérletazonosítóhoz tartozó tanulható embedding, gerjesztés-leíró paraméterek vagy külön hálózat futamonként. Önmagában a pillanatnyi feszültségérték nem mindig írja le az állapot múltját, így nem helyettesíti automatikusan a bemeneti trajektóriát és a kezdeti feltételeket.

Inverz feladatnál a kiválasztott fizikai paraméterek globális tanítható változóként kerülnek a modellbe. Ismeretlen terhelés esetén kiegészítő kimenet becsülheti $\hat{M}_L(t)$-t. A hálózat mélysége és szélessége hiperparaméter, nem fizikai adottság; kis, stabil modellről érdemes indulni, majd validáció alapján bővíteni.

\subsection{Fizikai maradékok}

Az elektromos, mechanikai és kinematikai egyenletek becsült maradékai:
\begin{align}
  r_e(t) &=
  u_a(t)-R\hat{i}(t)-L\dot{\hat{i}}(t)-k_e\hat{\omega}(t),
  \label{eq:re}\\
  r_m(t) &=
  J\dot{\hat{\omega}}(t)-k_t\hat{i}(t)
      +b\hat{\omega}(t)+\hat{M}_L(t),
  \label{eq:rm}\\
  r_{\theta}(t) &= \dot{\hat{\theta}}(t)-\hat{\omega}(t).
  \label{eq:rtheta}
\end{align}
Ha az elektromos mennyiségek gyors és a mechanikai mennyiségek lassú időskálán változnak, a kollokációs pontokat is megfelelően sűrűn kell elosztani; a lépésközt a mintavétel és a várható elektromos időállandó $\tau_e=L/R$ alapján kell megválasztani.

\subsection{Teljes veszteség}

Legyen $\mathcal{D}$ a mérések, $\mathcal{C}$ a kollokációs időpontok és $\mathcal{I}$ a kezdeti feltételek halmaza. A mérési veszteség csak a ténylegesen megfigyelt jelekre számítandó:
\begin{equation}
  \mathcal{L}_{\mathrm{data}} =
  \frac{1}{N_d}\sum_{j\in\mathcal{D}}\sum_{q\in\mathcal{Q}_j}
  w_q\left(
    \frac{\hat{y}_{q}(t_j)-y_{q,j}}{s_q}
  \right)^2,
  \label{eq:ldata}
\end{equation}
ahol $\mathcal{Q}_j$ a $t_j$ időpontban megfigyelt jelek halmaza, $s_q$ pedig a mértékegység- és nagyságrend-normalizáláshoz használt skála. A fizikai veszteség:
\begin{equation}
\begin{aligned}
  \mathcal{L}_{e} &=
  \frac{1}{N_c}\sum_{t_j\in\mathcal{C}}
    \left(\frac{r_e(t_j)}{s_e}\right)^2,\\
  \mathcal{L}_{m} &=
  \frac{1}{N_c}\sum_{t_j\in\mathcal{C}}
    \left(\frac{r_m(t_j)}{s_m}\right)^2,\\
  \mathcal{L}_{\theta} &=
  \frac{1}{N_c}\sum_{t_j\in\mathcal{C}}
    \left(\frac{r_{\theta}(t_j)}{s_{\dot{\theta}}}\right)^2,\\
  \mathcal{L}_{\mathrm{phys}} &= \mathcal{L}_e+\mathcal{L}_m+\mathcal{L}_{\theta}.
\end{aligned}
\label{eq:lphys}
\end{equation}
Az $s_e,s_m,s_{\dot{\theta}}$ maradékskálák rendre feszültség-, nyomaték- és szögsebesség-dimenziójúak. A normalizálás azért szükséges, mert a különböző fizikai egyenletek numerikus nagyságrendje eltérhet.

A kezdeti feltételek vesztesége például
\begin{equation}
  \mathcal{L}_{\mathrm{IC}}=
  \left(\frac{\hat{i}(t_0)-i_0}{s_i}\right)^2+
  \left(\frac{\hat{\omega}(t_0)-\omega_0}{s_\omega}\right)^2+
  \left(\frac{\hat{\theta}(t_0)-\theta_0}{s_\theta^{(0)}}\right)^2.
\end{equation}
A teljes célfüggvény:
\begin{equation}
\boxed{
  \mathcal{L}=
  \lambda_d\mathcal{L}_{\mathrm{data}}+
  \lambda_e\mathcal{L}_{e}+
  \lambda_m\mathcal{L}_{m}+
  \lambda_{\theta}\mathcal{L}_{\theta}+
  \lambda_0\mathcal{L}_{\mathrm{IC}}+
  \lambda_M\mathcal{L}_{\mathrm{load}}+
  \lambda_p\mathcal{L}_{\mathrm{prior}}.
}
\label{eq:loss}
\end{equation}
Az egyenletmaradékok külön jelölése az egyes tagok súlyának hangolását teszi lehetővé. $\mathcal{L}_{\mathrm{load}}$ csak akkor használatos, ha a terhelés időfüggvényére indokolt simasági vagy szakaszonkénti prior; $\mathcal{L}_{\mathrm{prior}}$ ismert adatlapértékek vagy fizikai korlátok gyenge priorját jelölheti. A $\lambda$ súlyokat és a skálákat validáció alapján kell rögzíteni. Zajos mérésnél négyzetes hiba helyett vagy mellett robusztus Huber-veszteség is vizsgálható.

\section{Ötlépcsős projektterv}

\begin{enumerate}[label=\textbf{\arabic*.},leftmargin=*]
  \item \textbf{Alapmodell és előreirányú feladat.}
  Ismert paraméterekkel tanítsuk meg a PINN-t az áram és a szögsebesség, majd a pozíció előrejelzésére. Először szimulált, teljes állapotadatot használjunk, és hasonlítsuk össze közvetlen numerikus ODE-megoldással.

  \item \textbf{Inverz feladat és paraméterazonosítás.}
  Ismert szimulációs gerjesztés mellett válasszunk ki kevés, azonosíthatónak tűnő paramétert (például $R,L$ vagy $J,b$), és becsüljük őket az állapotokkal együtt. Ezután ismételjük meg olyan kísérleten, ahol a paraméter referenciaértéke nem volt a tanítás része.

  \item \textbf{Ritka mérések és állapotbecslés.}
  A mérési pontok 10, 25, 50 és 100 százalékát felhasználva vizsgáljuk az áram-, sebesség- és pozícióbecslést. A pontok alulmintavételezését rögzített véletlen maggal vagy szabályos ritkítással tegyük reprodukálhatóvá.

  \item \textbf{Ismeretlen terhelés becslése.}
  A paramétereket először rögzítve rekonstruáljunk ismert, változó szimulációs terhelést. Csak ennek sikeres validálása után vizsgáljuk néhány paraméter és a terhelés együttes becslését.

  \item \textbf{Robusztusság és általánosítás.}
  Teszteljünk mérési zajt, ritkább mintavételezést, kisebb paramétereltérést, modellhibát és a tanításban nem látott bemeneti profilokat. Tartsunk külön, teljes futamokból álló teszthalmazt.
\end{enumerate}

Minden lépcsőben legyen dokumentált referencia, tanítóadat, validációs módszer és rögzített eredmény. A bonyolultabb feladatot csak akkor érdemes elkezdeni, ha az előző szint numerikusan és fizikailag ellenőrizhető eredményt ad.

\section{Validációs metrikák és összehasonlítás}

Minden kimenetre külön jelentsük a gyök-középnégyzetes hibát és az átlagos abszolút hibát:
\begin{align}
  \RMSE_y &= \sqrt{\frac{1}{N}\sum_{j=1}^{N}(\hat{y}_j-y_j)^2},\\
  \MAE_y &= \frac{1}{N}\sum_{j=1}^{N}|\hat{y}_j-y_j|.
\end{align}
Az áramra, sebességre, pozícióra és terhelőnyomatékra a mértékegységgel együtt kell közölni az eredményt. Kiegészítésként használható normalizált RMSE (a teszthalmaz szórásával vagy előre rögzített fizikai tartománnyal osztva), illetve $R^2$, de egyetlen összesített szám ne fedje el az egyes állapotok hibáját.

A fizikai konzisztencia méréséhez a tesztidőpontokban jelentsük az elektromos és mechanikai maradékok eredeti mértékegységű RMS-értékét. Paraméterazonosításnál jelentsük a relatív paraméterhibát:
\[
  \varepsilon_p =
  \frac{|\hat{p}-p_{\mathrm{ref}}|}{|p_{\mathrm{ref}}|}\,100\%.
\]
A terhelésbecsléshez külön $M_L$-MAE és -RMSE, ritka mérésnél pedig a megfigyelési arányt is közöljük.

Minimum összehasonlítások: (i) hagyományos adatvezérelt neurális hálózat; (ii) PINN sűrű méréssel; (iii) PINN ritka méréssel; (iv) numerikus fizikai modell, ha paraméterei ismertek. Azonos adatfelosztást és tesztfutamokat kell használni. A tanító-, validációs és teszthalmazt teljes kísérleti futam szerint célszerű szétválasztani, és az eredményt több tanítási magon is ellenőrizni.

\section{Fő kihívások és kockázatkezelés}

\begin{longtable}{>{\raggedright\arraybackslash}p{0.29\textwidth} >{\raggedright\arraybackslash}p{0.62\textwidth}}
\toprule
Kihívás & Kezelési mód \\
\midrule
\endhead
Paraméterek azonosíthatósága & Gerjesztés- és paraméterérzékenységi vizsgálat; becsült paraméterek számának korlátozása; ismert adatlapértékek rögzítése vagy prior használata.\\
PWM és kapocsfeszültség & A PWM parancsot ne azonosítsuk automatikusan $u_a$-val; ellenőrizzük a tápfeszültséget, a meghajtót és az elérhető feszültségmérést.\\
Eltérő időskálák & Normalizáljuk az időt és az állapotokat; ellenőrizzük a kollokációs pontok sűrűségét és a veszteségtagok nagyságrendjét.\\
Zajos deriváltak & Ne véges differenciával becsüljük közvetlenül a zajos fordulatszám deriváltját; használjuk a hálózat automatikus deriválását, és ellenőrizzük a fizikai maradékot.\\
Terhelés és súrlódás összemosódása & Előbb külön kalibráljuk a súrlódást, majd ismert szimulált terhelésen teszteljük a terhelésbecslést.\\
Szimuláció--mérés eltérés & A szintetikus és mért eredményt külön jelentsük; dokumentáljuk az érzékelőszűrést, késleltetést, telítést és az egyszerűsített modellből eredő hibát.\\
Túl optimista teszt & Időpontok véletlen keverése helyett teljes kísérleti futamokat tartsunk vissza; ne hangoljunk a teszthalmaz alapján.\\
\bottomrule
\end{longtable}

\section{Reprodukálhatóság és javasolt munkatermékek}

A projekt minden futása rögzítse a forrásadat verzióját és licencét, a fájlok ellenőrzött listáját, az egységkonverziókat, a bemeneti profilokat, a kezdeti feltételeket, a paramétereket, a mintavételezést, az adatfelosztást, a hálózati konfigurációt, a veszteségsúlyokat és a véletlen magot. A becsléseket eredeti fizikai mértékegységben is mentsük.

A végső projektcsomag minimális elemei:
\begin{itemize}
  \item MATLAB/Simulink referencia- és adatgeneráló modell;
  \item adatbetöltő és előfeldolgozó szkript egység- és jelnév-ellenőrzéssel;
  \item PINN tanító kód külön konfigurációval az öt lépcsőhöz;
  \item reprodukálható tesztkiértékelés és összehasonlító táblázatok;
  \item rövid adatlap a feltevésekről, sikertelen futásokról és korlátokról.
\end{itemize}

\begin{thebibliography}{9}
\bibitem{mendeley}
A. Von Chong és D. Cardenas,
\emph{Dataset for an end-to-end open-source DC motor control workflow: current calibration, system identification, and controller validation},
Mendeley Data, 2. verzió, 2026. szeptember 6.,
DOI: \href{https://doi.org/10.17632/5xvg43r9r8.2}{10.17632/5xvg43r9r8.2}.
Adatbázis-rekord: \url{https://data.mendeley.com/datasets/5xvg43r9r8/2}.
Hozzáférés dátuma: 2026. november 3.
\end{thebibliography}

\vfill
\begin{center}
\color{navy}\rule{0.65\textwidth}{0.5pt}\\
\textit{A dokumentum projektterv; a konkrét motorparaméterek, mérési jelek és üzemi korlátok kiválasztása külön adatellenőrzési lépést igényel.}
\end{center}

\end{document}
